\section{Reproducibility}
\label{app:repro}
All code, configurations and result tables are in the project repository.
Each run is defined by a YAML file in \texttt{configs/} (e.g.\ \texttt{padim\_wrn50.yaml}, \texttt{patchcore\_wrn50.yaml}) and executed with \texttt{python -m ad.run --config <file>}.
Tables are aggregated with \texttt{python -m ad.aggregate}, and the failure analysis is reproduced by \texttt{scripts/analysis\_failures.py} and \texttt{scripts/analysis\_failure\_panels.py}.
The tables and figures in this report are generated from the committed CSV files by \texttt{scripts/paper\_tables.py} and \texttt{scripts/paper\_figures.py}.
The protocol is enforced in code.
Models receive no labels or masks.
Thresholds and score normalisation are computed only from the normal validation split (\texttt{ad.threshold}).
Test labels are read only by \texttt{ad.metrics}.
Unit tests check these contracts.

\section{Recall by Defect Type}
\label{app:recall}
\begin{table}[h]
  \centering
  \caption{Image recall at the validation threshold, by defect type.}
  \label{tab:app_recall}
  \footnotesize
  \begin{tabular}{@{}l@{\hspace{4pt}}l@{\hspace{4pt}}rrr@{}}
\toprule
Category & Defect & $n$ & PaDiM & PatchCore \\
\midrule
bottle & broken\_large & 20 & \textbf{1.000} & \textbf{1.000} \\
bottle & broken\_small & 22 & \textbf{1.000} & \textbf{1.000} \\
bottle & contamination & 21 & \textbf{1.000} & \textbf{1.000} \\
\midrule
capsule & crack & 23 & \textbf{0.870} & 0.783 \\
capsule & faulty\_imprint & 22 & 0.455 & \textbf{0.682} \\
capsule & poke & 21 & 0.714 & \textbf{0.762} \\
capsule & scratch & 23 & 0.348 & \textbf{0.435} \\
capsule & squeeze & 20 & 0.950 & \textbf{1.000} \\
\midrule
screw & manipulated\_front & 24 & 0.083 & \textbf{0.667} \\
screw & scratch\_head & 24 & 0.083 & \textbf{0.917} \\
screw & scratch\_neck & 25 & 0.440 & \textbf{0.960} \\
screw & thread\_side & 23 & 0.000 & \textbf{0.522} \\
screw & thread\_top & 23 & 0.000 & \textbf{0.913} \\
\midrule
carpet & color & 19 & \textbf{1.000} & \textbf{1.000} \\
carpet & cut & 17 & \textbf{1.000} & \textbf{1.000} \\
carpet & hole & 17 & \textbf{1.000} & \textbf{1.000} \\
carpet & metal\_contamination & 17 & \textbf{1.000} & 0.941 \\
carpet & thread & 19 & \textbf{0.947} & \textbf{0.947} \\
\midrule
tile & crack & 17 & 0.353 & \textbf{1.000} \\
tile & glue\_strip & 18 & 0.167 & \textbf{1.000} \\
tile & gray\_stroke & 16 & 0.063 & \textbf{0.375} \\
tile & oil & 18 & 0.167 & \textbf{1.000} \\
tile & rough & 15 & 0.067 & \textbf{0.600} \\
\midrule
leather & color & 19 & 0.947 & \textbf{1.000} \\
leather & cut & 19 & \textbf{1.000} & \textbf{1.000} \\
leather & fold & 17 & \textbf{1.000} & \textbf{1.000} \\
leather & glue & 19 & \textbf{1.000} & \textbf{1.000} \\
leather & poke & 18 & \textbf{1.000} & \textbf{1.000} \\
\bottomrule
\end{tabular}

\end{table}

\section{Per-Category Results of the Enhancements}
\label{app:bonus}
\begin{table*}[h]
  \centering
  \caption{DINOv2 variants per category. \posthoc{}Post-hoc and test-informed.}
  \label{tab:app_dino}
  \footnotesize
  \begin{tabular}{@{}llrrr@{}}
\toprule
Category & Variant & \shortstack{Image\\AUROC} & \shortstack{Image F1\\(val thr)} & AU-PRO \\
\midrule
bottle & WRN50 baseline (Eq.~7) & 1.000 & 1.000 & 0.949 \\
bottle & DINOv2 ViT-B/14 & 1.000 & 0.992 & 0.962 \\
bottle & DINOv2 ViT-S/14 & 1.000 & 1.000 & 0.965 \\
bottle & WRN50 + DINOv2-B & 1.000 & 1.000 & 0.956 \\
bottle & DINOv2-B no-LN\posthoc & 1.000 & 1.000 & 0.926 \\
\midrule
capsule & WRN50 baseline (Eq.~7) & 0.978 & 0.836 & 0.931 \\
capsule & DINOv2 ViT-B/14 & 0.971 & 0.860 & 0.945 \\
capsule & DINOv2 ViT-S/14 & 0.990 & 0.853 & 0.940 \\
capsule & WRN50 + DINOv2-B & 0.979 & 0.802 & 0.940 \\
capsule & DINOv2-B no-LN\posthoc & 0.960 & 0.713 & 0.942 \\
\midrule
screw & WRN50 baseline (Eq.~7) & 0.970 & 0.884 & 0.963 \\
screw & DINOv2 ViT-B/14 & 0.879 & 0.674 & 0.845 \\
screw & DINOv2 ViT-S/14 & 0.822 & 0.738 & 0.836 \\
screw & WRN50 + DINOv2-B & 0.809 & 0.663 & 0.768 \\
screw & DINOv2-B no-LN\posthoc & 0.960 & 0.767 & 0.890 \\
\midrule
carpet & WRN50 baseline (Eq.~7) & 0.987 & 0.935 & 0.950 \\
carpet & DINOv2 ViT-B/14 & 1.000 & 0.978 & 0.980 \\
carpet & DINOv2 ViT-S/14 & 1.000 & 0.983 & 0.984 \\
carpet & WRN50 + DINOv2-B & 1.000 & 0.989 & 0.976 \\
carpet & DINOv2-B no-LN\posthoc & 0.997 & 0.962 & 0.944 \\
\midrule
tile & WRN50 baseline (Eq.~7) & 0.989 & 0.895 & 0.819 \\
tile & DINOv2 ViT-B/14 & 0.999 & 0.982 & 0.878 \\
tile & DINOv2 ViT-S/14 & 0.995 & 0.994 & 0.865 \\
tile & WRN50 + DINOv2-B & 0.999 & 0.982 & 0.870 \\
tile & DINOv2-B no-LN\posthoc & 1.000 & 1.000 & 0.800 \\
\midrule
leather & WRN50 baseline (Eq.~7) & 1.000 & 0.963 & 0.972 \\
leather & DINOv2 ViT-B/14 & 1.000 & 1.000 & 0.977 \\
leather & DINOv2 ViT-S/14 & 1.000 & 1.000 & 0.984 \\
leather & WRN50 + DINOv2-B & 1.000 & 1.000 & 0.975 \\
leather & DINOv2-B no-LN\posthoc & 1.000 & 0.989 & 0.945 \\
\bottomrule
\end{tabular}

\end{table*}

\begin{table*}[h]
  \centering
  \caption{Multi-scale variants per category.}
  \label{tab:app_ms}
  \footnotesize
  \begin{tabular}{@{}llrrr@{}}
\toprule
Category & Variant & \shortstack{Image\\AUROC} & \shortstack{Image F1\\(val thr)} & AU-PRO \\
\midrule
bottle & WRN50 baseline (Eq.~7) & 1.000 & 1.000 & 0.949 \\
bottle & Layers 1+2+3 & 1.000 & 0.984 & 0.962 \\
bottle & Layers 2+3+4 & 1.000 & 0.992 & 0.729 \\
bottle & Max-score control & 1.000 & 1.000 & 0.949 \\
bottle & Pyramid mean & 1.000 & 1.000 & 0.957 \\
bottle & $z$-pyramid mean & 1.000 & 1.000 & 0.957 \\
\midrule
capsule & WRN50 baseline (Eq.~7) & 0.978 & 0.836 & 0.931 \\
capsule & Layers 1+2+3 & 0.955 & 0.764 & 0.941 \\
capsule & Layers 2+3+4 & 0.895 & 0.893 & 0.725 \\
capsule & Max-score control & 0.974 & 0.866 & 0.931 \\
capsule & Pyramid mean & 0.974 & 0.848 & 0.939 \\
capsule & $z$-pyramid mean & 0.974 & 0.848 & 0.939 \\
\midrule
screw & WRN50 baseline (Eq.~7) & 0.970 & 0.884 & 0.963 \\
screw & Layers 1+2+3 & 0.978 & 0.895 & 0.982 \\
screw & Layers 2+3+4 & 0.665 & 0.065 & 0.633 \\
screw & Max-score control & 0.978 & 0.868 & 0.963 \\
screw & Pyramid mean & 0.984 & 0.929 & 0.972 \\
screw & $z$-pyramid mean & 0.984 & 0.929 & 0.972 \\
\midrule
carpet & WRN50 baseline (Eq.~7) & 0.987 & 0.935 & 0.950 \\
carpet & Layers 1+2+3 & 0.991 & 0.946 & 0.962 \\
carpet & Layers 2+3+4 & 0.997 & 0.967 & 0.752 \\
carpet & Max-score control & 0.987 & 0.937 & 0.950 \\
carpet & Pyramid mean & 0.993 & 0.937 & 0.961 \\
carpet & $z$-pyramid mean & 0.993 & 0.937 & 0.961 \\
\midrule
tile & WRN50 baseline (Eq.~7) & 0.989 & 0.895 & 0.819 \\
tile & Layers 1+2+3 & 0.996 & 0.817 & 0.844 \\
tile & Layers 2+3+4 & 0.978 & 0.976 & 0.624 \\
tile & Max-score control & 0.991 & 0.976 & 0.819 \\
tile & Pyramid mean & 0.991 & 0.976 & 0.833 \\
tile & $z$-pyramid mean & 0.991 & 0.976 & 0.833 \\
\midrule
leather & WRN50 baseline (Eq.~7) & 1.000 & 0.963 & 0.972 \\
leather & Layers 1+2+3 & 1.000 & 0.984 & 0.981 \\
leather & Layers 2+3+4 & 0.998 & 0.979 & 0.772 \\
leather & Max-score control & 1.000 & 0.958 & 0.972 \\
leather & Pyramid mean & 1.000 & 0.968 & 0.977 \\
leather & $z$-pyramid mean & 1.000 & 0.968 & 0.977 \\
\bottomrule
\end{tabular}

\end{table*}

\begin{table*}[h]
  \centering
  \caption{Post-processing variants per category.}
  \label{tab:app_pp}
  \footnotesize
  \begin{tabular}{@{}llrrr@{}}
\toprule
Category & Variant & \shortstack{Image\\AUROC} & \shortstack{Image F1\\(val thr)} & AU-PRO \\
\midrule
bottle & WRN50 baseline (Eq.~7) & 1.000 & 1.000 & 0.949 \\
bottle & Top-1\% mean & 1.000 & 1.000 & 0.949 \\
bottle & Smoothed max & 1.000 & 1.000 & 0.949 \\
bottle & Smoothing $\sigma=0$ & 1.000 & 1.000 & 0.951 \\
bottle & Smoothing $\sigma=2$ & 1.000 & 1.000 & 0.951 \\
bottle & Smoothing $\sigma=8$ & 1.000 & 1.000 & 0.942 \\
bottle & Global threshold $z=3$ & 1.000 & 1.000 & 0.949 \\
bottle & Global pooled val & 1.000 & 1.000 & 0.949 \\
\midrule
capsule & WRN50 baseline (Eq.~7) & 0.978 & 0.836 & 0.931 \\
capsule & Top-1\% mean & 0.976 & 0.922 & 0.931 \\
capsule & Smoothed max & 0.974 & 0.866 & 0.931 \\
capsule & Smoothing $\sigma=0$ & 0.978 & 0.836 & 0.933 \\
capsule & Smoothing $\sigma=2$ & 0.978 & 0.836 & 0.933 \\
capsule & Smoothing $\sigma=8$ & 0.978 & 0.836 & 0.921 \\
capsule & Global threshold $z=3$ & 0.978 & 0.710 & 0.931 \\
capsule & Global pooled val & 0.978 & 0.646 & 0.931 \\
\midrule
screw & WRN50 baseline (Eq.~7) & 0.970 & 0.884 & 0.963 \\
screw & Top-1\% mean & 0.975 & 0.829 & 0.963 \\
screw & Smoothed max & 0.978 & 0.868 & 0.963 \\
screw & Smoothing $\sigma=0$ & 0.970 & 0.884 & 0.964 \\
screw & Smoothing $\sigma=2$ & 0.970 & 0.884 & 0.964 \\
screw & Smoothing $\sigma=8$ & 0.970 & 0.884 & 0.955 \\
screw & Global threshold $z=3$ & 0.970 & 0.868 & 0.963 \\
screw & Global pooled val & 0.970 & 0.818 & 0.963 \\
\midrule
carpet & WRN50 baseline (Eq.~7) & 0.987 & 0.935 & 0.950 \\
carpet & Top-1\% mean & 0.983 & 0.937 & 0.950 \\
carpet & Smoothed max & 0.987 & 0.937 & 0.950 \\
carpet & Smoothing $\sigma=0$ & 0.987 & 0.935 & 0.952 \\
carpet & Smoothing $\sigma=2$ & 0.987 & 0.935 & 0.952 \\
carpet & Smoothing $\sigma=8$ & 0.987 & 0.935 & 0.942 \\
carpet & Global threshold $z=3$ & 0.987 & 0.935 & 0.950 \\
carpet & Global pooled val & 0.987 & 0.941 & 0.950 \\
\midrule
tile & WRN50 baseline (Eq.~7) & 0.989 & 0.895 & 0.819 \\
tile & Top-1\% mean & 0.992 & 0.982 & 0.819 \\
tile & Smoothed max & 0.991 & 0.976 & 0.819 \\
tile & Smoothing $\sigma=0$ & 0.989 & 0.895 & 0.819 \\
tile & Smoothing $\sigma=2$ & 0.989 & 0.895 & 0.819 \\
tile & Smoothing $\sigma=8$ & 0.989 & 0.895 & 0.816 \\
tile & Global threshold $z=3$ & 0.989 & 0.857 & 0.819 \\
tile & Global pooled val & 0.989 & 0.833 & 0.819 \\
\midrule
leather & WRN50 baseline (Eq.~7) & 1.000 & 0.963 & 0.972 \\
leather & Top-1\% mean & 1.000 & 0.939 & 0.972 \\
leather & Smoothed max & 1.000 & 0.958 & 0.972 \\
leather & Smoothing $\sigma=0$ & 1.000 & 0.963 & 0.975 \\
leather & Smoothing $\sigma=2$ & 1.000 & 0.963 & 0.974 \\
leather & Smoothing $\sigma=8$ & 1.000 & 0.963 & 0.965 \\
leather & Global threshold $z=3$ & 1.000 & 0.968 & 0.972 \\
leather & Global pooled val & 1.000 & 0.974 & 0.972 \\
\bottomrule
\end{tabular}

\end{table*}

\section{All Curated Failure Cases}
\label{app:failures}
Each panel shows the image with the GT outlined in green, the PaDiM map and the PatchCore map. Maps are divided by the validation pixel threshold, and the red contours mark the predicted regions.

\begin{figure*}[h]
  \centering
  \includegraphics[width=\textwidth]{fig_app_failures_1.jpg}
  \caption{Failure cases 1--6. (a)~Screw thread\_top: missed only by PaDiM (pose). (b)~Screw thread\_side: missed by both; two tiny regions on the normal high-frequency thread. (c)~Capsule scratch: missed by both; PaDiM's only hot spot is background dust. (d)~Capsule crack: missed by both; peaks on dark specks. (e)~Capsule faulty imprint: localised but not detected by PatchCore. (f)~Tile gray\_stroke: only about 10\% of the GT survives the centre crop.}
  \label{fig:app_fail1}
\end{figure*}

\begin{figure*}[h]
  \centering
  \includegraphics[width=\textwidth]{fig_app_failures_2.jpg}
  \caption{Failure cases 7--12. (a)~Tile oil: PaDiM calibration failure; PatchCore IoU 0.85. (b)~Normal carpet: global-shift false positive by PatchCore. (c)~Normal leather: diffuse PatchCore false positive, absorbed by PaDiM's looser threshold. (d)~Normal capsule: PaDiM-only false positive caused by background dust. (e)~Normal bottle: PaDiM-only false positive on the glossy rim. (f)~Leather colour: true positive with poor localisation (IoU 0.07).}
  \label{fig:app_fail2}
\end{figure*}

\section{Qualitative Results}
\label{app:qual}
Figs.~\ref{fig:app_qual_bottle}--\ref{fig:app_qual_leather} show the per-category qualitative summaries for PaDiM (left) and PatchCore (right).
Each method selects its own examples.
From top to bottom, the rows show true positives with the best localisation, true positives with poor localisation, missed anomalies (FN) and false alarms (FP).
The GT is outlined in green and the predicted region (map at or above the validation pixel threshold) in red.

\newcommand{\qualfig}[1]{%
\begin{figure*}[p]
  \centering
  \includegraphics[width=\textwidth,height=0.85\textheight,keepaspectratio]{fig_app_qual_#1.jpg}
  \caption{Qualitative results for \emph{#1}: PaDiM (left) and PatchCore (right).}
  \label{fig:app_qual_#1}
\end{figure*}}
\qualfig{bottle}
\qualfig{capsule}
\qualfig{screw}
\qualfig{carpet}
\qualfig{tile}
\qualfig{leather}
